\usepackage[T1]{fontenc}
\usepackage[utf8]{inputenc}
\usepackage{lmodern}
\usepackage[ruled,vlined,linesnumbered]{algorithm2e}
% ---- algorithm2e polish ----
\SetAlgoCaptionSeparator{.}          % 标题“Algorithm 1. 标题”
\SetKwComment{Comment}{$\triangleright$ }{}  % 行尾注释  \Comment{...}
\SetKwProg{Fn}{Function}{}{}         % \Fn{name(args)}{ ... }
\DontPrintSemicolon                  % 去掉行末分号（可选）
\SetAlCapFnt{\bfseries}              % caption 加粗
\SetKwInOut{Input}{Input}\SetKwInOut{Output}{Output} % 英文 In/Out
% =========================
% Page & Spacing (A4)
% =========================
\usepackage[a4paper,margin=1in,headsep=0.35in]{geometry}
\setlength{\parindent}{0pt}
\setlength{\parskip}{6pt}

% =========================
% Colors (aligned to your palette)
% =========================
\usepackage{xcolor}
\definecolor{brand}{HTML}{1A9D8F}      % teal
\definecolor{brandD}{HTML}{2A7F6F}     % dark teal
\definecolor{ink}{HTML}{1A1A1A}        % text
\definecolor{keybg}{HTML}{E8F4F1}      % light teal background
\definecolor{keydark}{HTML}{2F5E58}    % dark teal
\definecolor{accent}{HTML}{FF6B6B}     % coral
\definecolor{accentD}{HTML}{D94F4F}    % dark coral
\definecolor{soft}{HTML}{F3F8F7}
\definecolor{secblue}{HTML}{2B44C6}    % deeper blue for sections
\definecolor{warn}{HTML}{FFD54F}     % 黄色警告
\definecolor{info}{HTML}{64B5F6}     % 浅蓝信息
\definecolor{termblue}{HTML}{1A4C9A} % 深蓝（专业名词）
\definecolor{error}{HTML}{E53935}    % 红色错误
\color{ink}

% =========================
% Graphics / TikZ
% =========================
\usepackage{tikz,graphicx}
\usetikzlibrary{calc,positioning,arrows.meta,decorations.pathmorphing}

% =========================
% Hyperref
% =========================
\usepackage[
  colorlinks=true,
  linkcolor=secblue,
  citecolor=brand,
  urlcolor=brandD,
  bookmarksopen=true,
  bookmarksnumbered=true
]{hyperref}
\hypersetup{
  pdftitle    = {Your Note Title},
  pdfauthor   = {Your Name},
  pdfsubject  = {Academic Notes},
  pdfcreator  = {LaTeX with hyperref},
  pdfkeywords = {probability, statistics, optimization, learning theory}
}
\usepackage[final]{microtype}
\linespread{1.04} % or 1.05

% =========================
% Headers / Footers (every page)
%   - Right header: italicized current section title
%   - Footer: center course info; page number at right
%   - Chapter openings also use same style (redefine 'plain')
% =========================
\usepackage{fancyhdr}
\pagestyle{fancy}
\fancyhf{}
\renewcommand{\headrulewidth}{0.4pt}
\renewcommand{\sectionmark}[1]{\markright{\thesection\ \emph{#1}}}
\fancyhead[R]{\nouppercase{\rightmark}}
% 奇偶页页脚对称：左/右页页码与信息互换
\fancyfoot[LE,RO]{\textcolor{brand}{\thepage}}
\fancyfoot[LO,RE]{\textcolor{brand}{Notes}}
\fancypagestyle{plain}{%
  \fancyhf{}%
  \renewcommand{\headrulewidth}{0pt}% ← 没有页眉线
  % 不设 \fancyhead[...] → 没有页眉
  % 仍然保留你想要的页脚：
  \fancyfoot[LE,RO]{\textcolor{brand}{\thepage}}
  \fancyfoot[LO,RE]{\textcolor{brand}{Notes}}
}

% =========================
% Chapter & Section formatting
% =========================
% =========================
% Chapter & Section formatting
% =========================
\usepackage{titlesec}
% ===== Part 样式（居中 + 下划线；不影响 Chapter/Section）=====
\titleformat{\part}[display]
  {\filcenter\bfseries\color{keydark}}   % 颜色只在标题块内生效
  {\Large Part \thepart}                  % 标签行（Part I）
  {0.6em}                                 % 标签与标题间距
  {\Huge}                                 % 标题字号
  [\vspace{0.8em}{\color{brand}\rule{7cm}{1.2pt}}] % 结尾横线（局部上色）

% 整块垂直居中；若想稍微靠上，改成 {0.6\fill}{0.4\fill}
\titlespacing*{\part}{0pt}{\fill}{\fill}
% 自定义 Chapter 样式
\titleformat{\chapter}[display] % display 代表上下分行
  {\bfseries\Huge}              % 字体设置
  {Chapter \thechapter}         % 第一行显示的内容
  {0pt}                         % 两行之间的间距
  {\Huge}                       % 标题字体      % 在标题下加一条横线（可去掉）

\titleformat{\section}
  {\Large\bfseries\color{secblue}}
  {\thesection}{0.6em}{}
\titlespacing*{\section}{0pt}{0.9em}{0.45em}

\titleformat{\subsection}
  {\large\bfseries\color{secblue}}
  {\thesubsection}{0.6em}{}
\titlespacing*{\subsection}{0pt}{0.8em}{0.35em}
% ---- toggle chapter heading visibility ----
\newcommand{\HideChapterHeads}{%
  \titleformat{\chapter}[display]{\relax}{}{0pt}{}%
  \titlespacing*{\chapter}{0pt}{0pt}{0pt}%
}
\newcommand{\ShowChapterHeads}{%
  \titleformat{\chapter}
    {\huge\bfseries\color{black}}{Chapter \thechapter}{0.8em}{}%
  \titlespacing*{\chapter}{0pt}{-0.2em}{0.9em}%
}
% =========================
% SideBar block (kept)
% =========================
\usepackage{mdframed}
\newmdenv[
  topline=false, bottomline=false, rightline=false,
  linewidth=4pt, linecolor=brand, backgroundcolor=soft,
  innerleftmargin=10pt, innerrightmargin=10pt,
  skipabove=6pt, skipbelow=6pt
]{SideBar}

% =========================
% Shaded environments (framed)
%  - Theorem, Lemma, Definition: shaded + italic body
%  - Single shared counter across all theorem-like envs
% =========================
\usepackage{framed}
\usepackage{enumitem}
\definecolor{shadecolor}{RGB}{230,240,250} % light blue shaded background

\usepackage{amsthm,amsmath,amssymb,bm}
\usepackage{cleveref} % 智能引用，必须在 hyperref 之后
\numberwithin{equation}{section}
\DeclareMathOperator{\Var}{Var}
\newcommand{\E}{\mathbb{E}}

% master counter per section
\newtheorem{thm}{Theorem}[section]
\newtheorem{lem}[thm]{Lemma}
\newtheorem{defn}[thm]{Definition}
\newtheorem{prop}[thm]{Proposition}
\newtheorem{cor}[thm]{Corollary}
\newtheorem{example}[thm]{Example}
\newtheorem{que}[thm]{Question}
\newtheorem{assumption}[thm]{Assumption}
\newtheorem{exrc}[thm]{Exercise}



% Shaded wrappers (italic body)
\newenvironment{theoremS}[1][]{%
  \begin{shaded}%
  \if\relax\detokenize{#1}\relax
    \begin{thm}\itshape
  \else
    \begin{thm}[\textbf{#1}]\itshape
  \fi
}{%
  \end{thm}%
  \end{shaded}%
}
\newenvironment{lemmaS}[1][]{%
  \begin{shaded}\begin{lem}[#1]\itshape}{\end{lem}\end{shaded}}
\newenvironment{definitionS}[1][]{%
  \begin{shaded}\begin{defn}[#1]\itshape}{\end{defn}\end{shaded}}

% 让 exercise 使用 theorem 计数器


% 让 exercise 使用 theorem 计数器
\newenvironment{exercise}[1][]{%
  \par\smallskip
  \refstepcounter{thm}% 推进 theorem 计数器
  \begingroup\bfseries\color{keydark}Exercise \thethm%
  \if\relax\detokenize{#1}\relax\else\ \ (\textit{#1})\fi
  \par\medskip\endgroup
}{\par\smallskip}

% RemarkNote: red title, part of same counter, non-shaded
\newenvironment{RemarkNote}[1][]{%
  \par\smallskip
  \noindent{\bfseries\color{accentD}Remark}%
  \ifx\relax#1\relax\else{\bfseries\color{accentD} \ (#1)}\fi
  \ .\ %
}{%
  \par\smallskip
  {\color{ink!60}\rule{\linewidth}{0.4pt}}
  \smallskip
}


% =========================
% Book Master Cover (for the whole book) — redesigned
% (uses your palette + math watermarks; consistent, elegant)
% =========================
\newcommand{\CoverYOffset}{0.00\paperheight}
\newcommand{\CoverCardWidth}{0.78\paperwidth}
\newcommand{\CoverTitleSize}{\fontsize{32pt}{36pt}\selectfont}
\newcommand{\CoverSubtitleSize}{\fontsize{16pt}{20pt}\selectfont}
\newcommand{\CoverMetaSize}{\fontsize{11pt}{14pt}\selectfont}
\newcommand{\CoverSmall}{\fontsize{10pt}{13pt}\selectfont}
\newcommand{\BookDate}{\today}

\newcommand{\MakeMasterBookCover}[7]{%
  \thispagestyle{empty}%
  \begin{tikzpicture}[remember picture, overlay]

    % Background halves
    \fill[keybg] (current page.north west) rectangle (current page.south east);
    \fill[white] ($(current page.north west)+(0.48\paperwidth,0)$)
                 rectangle (current page.south east);

    % Very light geometric grid on the right
    \foreach \x in {0.52,0.58,0.64,0.70,0.76,0.82,0.88,0.94} {
      \draw[brand!12, line width=0.3pt]
        ($(current page.north west)+(\x\paperwidth,0.8cm)$) --
        ($(current page.south west)+(\x\paperwidth,-0.8cm)$);
    }
    \foreach \y in {0.18,0.32,0.46,0.60,0.74,0.88} {
      \draw[brand!10, line width=0.3pt]
        ($(current page.west)+(0.55\paperwidth,\y\paperheight)$) --
        ($(current page.east)+(-0.55\paperwidth,\y\paperheight)$);
    }

    % Soft circles (math-ish motif)
    \foreach \r/\x/\y in {1.4/2.0/2.8, 0.9/5.8/8.5, 1.2/11.5/5.2}{
      \draw[keydark!10, line width=8pt]
        ($(current page.south west)+(\x,\y)$) circle (\r);
    }

    % Low-opacity math watermarks
    \node[rotate=14, scale=0.95, text opacity=0.07]
      at ($(current page.center)+(-5.6cm,5.0cm)$)
      {$\displaystyle \mathbb{E}\!\left[e^{\lambda X}\right] \le \exp\!\Big(\tfrac{\lambda^2 \sigma^2}{2}\Big)$};

    \node[rotate=-18, scale=0.9, text opacity=0.065]
      at ($(current page.center)+(6.3cm,-5.4cm)$)
      {$\displaystyle \Pr\!\big(|X^\top A X - \mathbb{E}[X^\top A X]|\ge t\big)
        \le 2\exp\!\Big(-c\,\min\{\tfrac{t^2}{K^4\|A\|_F^2},\tfrac{t}{K^2\|A\|}\}\Big)$};

    \node[rotate=28, scale=0.95, text opacity=0.06]
      at ($(current page.center)+(7.0cm,6.0cm)$)
      {$\displaystyle \mathrm{KL}(P\|Q)=\int \log\!\frac{\mathrm{d}P}{\mathrm{d}Q}\,\mathrm{d}P$};

    % Central Card
    \node[
      anchor=center,
      fill=white, draw=brandD!30, rounded corners=9pt,
      inner sep=16pt
    ] at ($ (current page.center) + (0,\CoverYOffset) $)
    {%
      \begin{minipage}{\CoverCardWidth}
        \raggedright

        % Tag (left) + seal (right)
        \noindent
        \tikz[baseline]{
          \node[
            fill=brand, rounded corners=3pt,
            inner xsep=8pt, inner ysep=4pt,
            text=white, font=\bfseries\small
          ] {Concentration Inequalities \& High-Dimensional Probability};
        }%
        \hfill
        \tikz[baseline]{
          \node[
            draw=brandD!60, rounded corners=3pt,
            inner xsep=8pt, inner ysep=3pt,
            text=brandD!90, font=\bfseries\scriptsize, fill=keybg
          ] {Graduate Notes};
        }\par

        % Title
        \vspace{0.7em}
        {{\CoverTitleSize\bfseries\color{ink} #1\par}}

        % Divider
        \vspace{0.35em}
        {\color{brandD!45}\rule[0.6ex]{7.2cm}{0.6pt}}\par

        % Subtitle
        \vspace{0.25em}
        {{\CoverSubtitleSize\color{secblue} #2\par}}

        % Author + date
        \vspace{0.8em}
        {\CoverMetaSize
          \textcolor{ink!85}{\textbf{#3}} \quad
          \textcolor{ink!55}{\(\cdot\)} \quad
          \textcolor{ink!70}{\BookDate}
          \par
        }

        % Purpose
        \vspace{0.9em}
        {\CoverMetaSize\textbf{\color{keydark}Purpose.}}
        {\CoverMetaSize\color{ink!90} #5\par}

        % What you will learn
        \vspace{0.7em}
        {\CoverMetaSize\bfseries\color{keydark}What you will learn}\par
        \vspace{0.1em}
        {\CoverSmall\color{ink!92}%
          #6
        }

        % What problems this book solves
        \vspace{0.6em}
        {\CoverMetaSize\bfseries\color{keydark}What problems this book solves}\par
        \vspace{0.1em}
        {\CoverSmall\color{ink!92}%
          #7
        }

      \end{minipage}
    };

    % Subtle footer right
    \node[anchor=south east,align=right,text=keydark!70,font=\small]
      at ([xshift=-0.6cm,yshift=0.7cm]current page.south east)
      {Compiled on \BookDate};

  \end{tikzpicture}
  \clearpage
}

% =========================
% Chapter Landing (per-chapter mini cover) — kept
% =========================
% Intuition / Analysis Blocks (English, polished; no varwidth)
% ===== Dependencies =====
\usepackage[most]{tcolorbox}
\usepackage{xcolor}
\usepackage{tikz}

% ===== Color palette (example; keep yours if already defined) =====
% \definecolor{brand}{HTML}{1A9D8F}
% \definecolor{brandD}{HTML}{2A7F6F}
% \definecolor{secblue}{HTML}{1A4F9D}
% \definecolor{ink}{gray}{0.1}

% ===== Global tcolorbox defaults =====
\tcbset{
  enhanced,
  breakable,
  sharp corners=all,
  boxsep=6pt,
  boxrule=0pt,          % no outer frame line
  colback=white,        % white background
  colframe=white,       % frame color = background, so invisible
  fonttitle=\bfseries\small,
  borderline west={0pt}{0pt}{white}, % default: no left bar unless overridden
}

% ===== helper: heading line (manual title) =====
\newcommand{\BoxHeading}[2]{%
  {\bfseries\small\textcolor{#1}{#2}}\par\medskip
}

% ===== Intuition box =====
\newtcolorbox{Intuition}[1][]{%
  colback=white,
  colframe=white, % ensure no frame except custom left bar
  boxrule=0pt,
  borderline west={2pt}{0pt}{brand!80}, % ONLY left bar
  before upper=\BoxHeading{brand!80}{Intuition},
  #1
}

% ===== Analysis box =====
\newtcolorbox{Analysis}[1][]{%
  colback=white,
  colframe=white,
  boxrule=0pt,
  borderline west={2pt}{0pt}{secblue!80},
  before upper=\BoxHeading{secblue!80}{Analysis},
  #1
}


% ===== Inline callouts =====
\newcommand{\intuition}[1]{%
  {\begingroup
   \tikz[baseline=-0.6ex]\draw[brand!80, line width=2pt] (0,0) -- (0,1.1ex);~%
   \textcolor{brand!80}{\footnotesize\textsf{#1}}%
   \endgroup}%
}

\newcommand{\analysis}[1]{%
  {\begingroup
   \tikz[baseline=-0.6ex]\draw[secblue!90, line width=2pt] (0,0) -- (0,1.1ex);~%
   \textcolor{secblue!80}{\footnotesize\textsf{#1}}%
   \endgroup}%
}
% =========================
% ====== Elegant Question Box (no blur keys; stable) ======
% 独立计数器：按 section 编号  <section>.<number>
% ===== Journal-style Problem (print-friendly, minimal) =====
% 编号：按 section 计数  <section>.<num>

\newcounter{prob}[section]
\renewcommand{\theprob}{\thesection.\arabic{prob}}

% 颜色：默认复用你的 secblue；如需黑白打印，见“单色开关”
\colorlet{pframe}{secblue!70}
\colorlet{pribbon}{secblue!90}
\colorlet{ptitle}{secblue!95}

% 单色/彩色切换（可选）：把 \def\PROBLEM_MONO{1} 打开即走黑白
% \def\PROBLEM_MONO{1}
\makeatletter
\@ifundefined{PROBLEM_MONO}{}{%
  \colorlet{pframe}{black!55}
  \colorlet{pribbon}{black}
  \colorlet{ptitle}{black}
}
\makeatother

% 样式：细边框 + 左侧 2pt 彩带 + 小型大写标题，无阴影/渐变
\newtcolorbox{ProblemBox}[1][]{%
  enhanced, breakable,
  colback=white, colframe=pframe, boxrule=0.5pt,
  borderline west = {2pt}{0pt}{pribbon},
  arc=3pt, outer arc=3pt,
  left=9pt, right=9pt, top=8pt, bottom=8pt,
  colbacktitle=white, coltitle=ptitle,
  fonttitle=\bfseries\scshape\small,  % 小型大写，期刊感
  toptitle=1pt, bottomtitle=1pt,
  title={Problem~\theprob%
    \if\relax\detokenize{#1}\relax\else\ \textnormal{\footnotesize[\,#1\,]}\fi},
  attach boxed title to top left = {xshift=2mm, yshift=-2mm}
}

% 封装环境：\begin{Problem}[可选标签] ... \end{Problem}
\newenvironment{Problem}[1][]{%
  \refstepcounter{prob}%
  \begin{ProblemBox}[#1]%
}{%
  \end{ProblemBox}%
}

% cleveref 名称映射
\crefname{prob}{Problem}{Problems}
\Crefname{prob}{Problem}{Problems}
\providecommand{\probautorefname}{Problem}
% ===========================================================

% =========================
% Technique Summary (red style, same as Analysis)
% =========================
\newtcolorbox{Technique}[1][]{%
  colback=white, colframe=accentD!70,
  colbacktitle=accentD,
  title={\textbf{Technique}},
  fonttitle=\bfseries\small,
  attach boxed title to top left = {yshift=-2mm, xshift=2mm},
  borderline west = {2pt}{0pt}{accentD},
  breakable,
  #1
}
% ===== Chapter cover knobs (尺寸等可按需调) =====

\usepackage[
  backend=biber,
  style=numeric,
  maxbibnames=9,
  maxcitenames=2,
  doi=false,
  isbn=false,
  url=true
]{biblatex}
\DefineBibliographyStrings{english}{%
  urlseen = {accessed},
}
\AtEveryBibitem{\clearfield{month}\clearfield{day}}
\addbibresource{refs.bib}
\addbibresource{references.bib}
% ==== Highlight System ====
\usepackage{xcolor,tikz}
\usepackage{soul} % 提供 \hl


% light backgrounds for theorem-like boxes
\definecolor{theoremBg}{RGB}{255,240,240} % 很浅的粉/暖红底；不刺眼
\colorlet{theoremBar}{accentD}            % 左侧竖条 = 你定义过的深珊瑚红 accentD

\colorlet{definitionBg}{keybg}            % 定义用你现成的浅绿 keybg
\colorlet{definitionBar}{brand}           % 左条 = brand (teal)

\definecolor{exampleBg}{RGB}{252,242,230} % 很浅米橘
\colorlet{exampleBar}{accentD}            % 左条 = accentD (深珊瑚)

\newcommand{\hltext}[2][yellow!30]{%
  \tikz[baseline=(X.base)] \node[rectangle, rounded corners=2pt,
  fill=#1, inner sep=2pt, anchor=base] (X) {#2};%
}

% ==== Semantic Shortcuts ====
\newcommand{\KeyIdea}[1]{\hltext[keybg]{\textbf{#1}}}            % 关键想法：浅绿底
\newcommand{\Term}[1]{\textcolor{termblue}{\textbf{#1}}}           % 专业名词：品牌深绿字
\newcommand{\Emph}[1]{\textcolor{accentD}{\textbf{#1}}}          % 强调词汇：红
\newcommand{\Confuse}[1]{\hltext[error!20]{\textcolor{error}{#1}}} % 不懂/存疑处

\newcommand{\Info}[1]{\textcolor{info!90!black}{#1}}             % 说明或注解
%一侧笔记


% 轻量 todo（不依赖 todonotes 包）
%\newcommand{\todo}[1]{%
%  \ifreview \hltext[warn!35]{\textbf{TODO:} #1}\else\textcolor{warn!60!black}{[TODO: #1]}\fi}
% ===== mdframed-wrapped theorem-like environments =====
% Generic geometry knobs
\newcommand{\TheoremFrameSetup}[2]{%
  \begin{mdframed}[
    topline=false,
    bottomline=false,
    rightline=false,
    leftline=true,
    linewidth=1.5pt,          % 左侧竖条宽度，可调 1pt-2pt
    linecolor=#1,             % 竖条颜色
    backgroundcolor=#2,       % 背景色
    innerleftmargin=10pt,
    innerrightmargin=10pt,
    innertopmargin=8pt,
    innerbottommargin=8pt,
    skipabove=10pt,
    skipbelow=10pt,
    roundcorner=2pt
  ]%
}
\newcommand{\EndTheoremFrame}{%
  \end{mdframed}%
}

\let\oldthm\thm
\let\endoldthm\endthm
\renewenvironment{thm}[1][]{%
  \TheoremFrameSetup{termblue}{shadecolor}%
  \if\relax\detokenize{#1}\relax
    \begin{oldthm}\itshape
  \else
    \begin{oldthm}[\textbf{#1}]\itshape
  \fi
}{%
  \end{oldthm}%
  \EndTheoremFrame
}

% --- Lemma (shares same counter as thm, already) ---
\let\oldlem\lem
\let\endoldlem\endlem
\renewenvironment{lem}[1][]{%
  \TheoremFrameSetup{theoremBar}{theoremBg}%
  \if\relax\detokenize{#1}\relax
    \begin{oldlem}\itshape
  \else
    \begin{oldlem}[\textbf{#1}]\itshape
  \fi
}{%
  \end{oldlem}%
  \EndTheoremFrame
}

% --- Proposition ---
\let\oldprop\prop
\let\endoldprop\endprop
\renewenvironment{prop}[1][]{%
  \TheoremFrameSetup{theoremBar}{theoremBg}%
  \if\relax\detokenize{#1}\relax
    \begin{oldprop}\itshape
  \else
    \begin{oldprop}[\textbf{#1}]\itshape
  \fi
}{%
  \end{oldprop}%
  \EndTheoremFrame
}

% --- Definition (green theme) ---
\let\olddefn\defn
\let\endolddefn\enddefn
\renewenvironment{defn}[1][]{%
  \TheoremFrameSetup{definitionBar}{definitionBg}%
  \if\relax\detokenize{#1}\relax
    \begin{olddefn}\itshape
  \else
    \begin{olddefn}[\textbf{#1}]\itshape
  \fi
}{%
  \end{olddefn}%
  \EndTheoremFrame
}

% --- Example (orange theme) ---
\let\oldexample\example
\let\endoldexample\endexample
\renewenvironment{example}[1][]{%
  \TheoremFrameSetup{exampleBar}{exampleBg}%
  \if\relax\detokenize{#1}\relax
    \begin{oldexample}\itshape
  \else
    \begin{oldexample}[\textbf{#1}]\itshape
  \fi
}{%
  \end{oldexample}%
  \EndTheoremFrame
}

\definecolor{QuestionBright}{HTML}{0076D6} % 更亮的蓝色，醒目不刺眼



\renewenvironment{que}[1][]{%
  \par\vspace{0.8em}%
  \noindent
  % 左侧彩色竖条
  \begingroup
    \color{QuestionBright}\rule[2pt]{2pt}{1.2em}%
  \endgroup
  \hspace{0.6em}%
  % 主标题行
  \textbf{\textcolor{QuestionBright}{Question~\thethm}}%
  \if\relax\detokenize{#1}\relax
    \textbf{\textcolor{QuestionBright}{.}}%
  \else
    \textbf{\textcolor{QuestionBright}{~(#1).}}%
  \fi
  \par\vspace{0.4em}%
  % 这里开始包一个 L 形的小框：左线 + 很短的底线
  \begin{mdframed}[
    topline=false,
    rightline=false,
    leftline=true,
    bottomline=true,
    linecolor=QuestionBright,
    linewidth=0.6pt,
    innerleftmargin=6pt,
    innerrightmargin=0pt,
    innertopmargin=2pt,
    innerbottommargin=2pt,
    skipabove=0pt,
    skipbelow=0pt,
    leftmargin=0pt,
    rightmargin=0.78\linewidth % 调这个数→底线就会很短，只露一点
  ]
  \begin{minipage}[t]{0.93\linewidth}%
  \setlength{\parindent}{1.2em}%
  \normalfont\itshape
}{%
  \end{minipage}%
  \end{mdframed}%
  \par\vspace{0.8em}%
}
% =========================================
% Custom proof environment (no italics body)
% Appearance: left gray bar, white bg, grey "Proof."
% =========================================

% 备份 amsthm 的原始 proof 环境，避免丢功能（比如 \qedhere）

\let\OldProof\proof
\let\EndOldProof\endproof

% 我们重写 proof 环境：
\usepackage{amsthm}          % 提供 \qedsymbol

\renewenvironment{proof}[1][]{%
  \begin{tcolorbox}[
    enhanced,
    breakable,
    colback=white,
    boxrule=0pt,
    borderline west={2pt}{0pt}{black!25},
    left=10pt,
    right=10pt,
    top=6pt,
    bottom=6pt,
    sharp corners,
    before skip=10pt,
    after skip=10pt
  ]%
  {\small\bfseries\color{black!55}Proof%
    \if\relax\detokenize{#1}\relax
    \else\ \textnormal{(#1)}%
    \fi.\par}%
  \vspace{4pt}%
  \normalfont
}{%
  \hfill\qedsymbol
  \end{tcolorbox}
}

% 自定义 QED 符号的样式（淡灰方块）
\renewcommand{\qedsymbol}{\textcolor{black!40}{$\blacksquare$}}


% 自定义 QED 符号的样式（淡灰方块）
\renewcommand{\qedsymbol}{\textcolor{black!40}{$\blacksquare$}}

\usepackage{float}  % 支持 [H] 参数

\usepackage{booktabs}
\usepackage{multirow}
% =======================================
% Common mathematical macros for probability, stochastic processes, and optimization.
% =======================================

% requires amsmath, amssymb, bbm
\usepackage{bbm}

% ----- Probability -----
\providecommand{\E}{\mathbb{E}}
\newcommand{\EE}[1]{\mathbb{E}\left[#1\right]}
\renewcommand{\P}{\mathbb{P}}
\providecommand{\Var}{\mathrm{Var}}
\newcommand{\Cov}{\mathrm{Cov}}
\newcommand{\KL}{\mathrm{KL}}
\newcommand{\TV}{\mathrm{TV}}
\newcommand{\1}{\mathbbm{1}}

% ----- Stochastic processes / SDE -----
\newcommand{\F}{\mathcal{F}}
\newcommand{\Ft}{\mathcal{F}_t}
\newcommand{\W}{W_t}
\newcommand{\dW}{\,dW_t}

% ----- Optimization -----
\newcommand{\grad}{\nabla}
\newcommand{\hess}{\nabla^2}
\newcommand{\ip}[1]{\left\langle #1 \right\rangle}
\newcommand{\norm}[1]{\left\lVert #1 \right\rVert}
\newcommand{\Div}{\mathrm{div}}
\newcommand{\Lap}{\Delta}
\DeclareMathOperator*{\argmin}{arg\,min}
\DeclareMathOperator*{\argmax}{arg\,max}

% ----- Diffusion / Score -----
\newcommand{\score}{\nabla_x \log p}
\newcommand{\sphi}{s_\phi}
\newcommand{\N}{\mathcal{N}}


% ----- Sets & Linear Algebra -----
\newcommand{\R}{\mathbb{R}}
\newcommand{\Rd}{\mathbb{R}^d}
\newcommand{\T}{^\top}
% ============================
% RL (Reinforcement Learning) macros
% ============================

% MDP
\newcommand{\Ssp}{\mathcal{S}}
\newcommand{\Asp}{\mathcal{A}}
\newcommand{\Pt}{\mathbb{P}}
\newcommand{\Psa}{P(\cdot\mid s,a)}


% Value / Q / A
\newcommand{\Vpi}{V^\pi}
\newcommand{\Qpi}{Q^\pi}
\newcommand{\Api}{A^\pi}

% Bellman operators
\newcommand{\Tpi}{\mathcal{T}^\pi}
\newcommand{\Tstar}{\mathcal{T}^\star}

% policy
\newcommand{\pisa}{\pi(a\mid s)}
\newcommand{\Pispace}{\Pi}
\newcommand{\pib}{\pi_{\text{b}}}
\newcommand{\pitr}{\pi_{\text{t}}}


% policy gradient / actor-critic
\newcommand{\gradJ}{\nabla J(\pi)}
\newcommand{\logpi}{\log \pi(a\mid s)}
\newcommand{\thetaParam}{\theta}

% regret
\newcommand{\Reg}{\mathrm{Reg}}
\newcommand{\RegT}{\mathrm{Reg}(T)}
